\section{Introduction}\label{sec:introduction}

The open-ended lab of the Software Construction course asks each group of two students to build a working task
management (TODO) system in Java, and to build it with the tools software is constructed with in practice: an IDE and
the JDK to write, run and debug the program, Git and GitHub to record its revisions, and the Javadoc tool to document
it. The program and the record of how it was made are assessed together.

\subsection{Objective}

The lab brief states two objectives, each tied to a course learning outcome (CLO):
\begin{itemize}
  \item \textbf{CLO-5:} independently build a working Java task management system, with recorded Git revisions and
        generated Javadoc documentation;
  \item \textbf{CLO-4:} practise and adapt modern tools (IDE, JDK, Git, GitHub and Javadoc) for software construction.
\end{itemize}
Table~\ref{tab:clo} shows how the two outcomes map to the programme learning outcomes and how much each one weighs.
Our own aim was to meet every functional requirement with input validation and clean error handling, and to leave
evidence for each step of the workflow (the IDE, the Git history, the Javadoc HTML and the test results) that can be
checked against the repository.

\begin{table}[!t]
  \caption{CLO mapping of the lab}\label{tab:clo}
  \centering\footnotesize
  \begin{tabularx}{\linewidth}{@{}lllXr@{}}
    \toprule
    CLO   & PLO   & Level & Learning outcome & Weight \\
    \midrule
    CLO-4 & PLO-5 & P-3   & Practise and adapt modern tools for software construction. & 25\% \\
    CLO-5 & PLO-5 & P-4   & Execute a software construction workflow independently using development,
                            version-control and documentation tools, producing a working application with recorded
                            revisions and generated code documentation. & 75\% \\
    \bottomrule
  \end{tabularx}
\end{table}

\subsection{Organisation of the report}

Section~\ref{sec:problem} states the problem and the requirements, and Section~\ref{sec:tools} the tools and the
environment. Section~\ref{sec:design} explains the design of the program and the reasons for it, and
Section~\ref{sec:dsa} its data structures and algorithms. Section~\ref{sec:functionality} walks through the program's
functions with screenshots. Sections~\ref{sec:ide}, \ref{sec:git} and~\ref{sec:javadoc} describe how the IDE, Git and
Javadoc were used, Section~\ref{sec:testing} gives the test cases with their expected and actual results, and
Section~\ref{sec:conclusion} concludes.
